\chapter{Vaginal Prolapse (Pelvic Organ Prolapse)}

\section*{Definition}
Vaginal dropping of an organ is the descent of the vaginal walls or uterus into the vaginal canal. It occurs when the pelvic floor muscles and connective tissue are weakened or damaged. Types include cystocele (bladder), rectocele (rectum), and uterine dropping of an organ.

\section*{Pathophysiology}
The pelvic floor is a hammock-like structure of muscles and fascia that supports the pelvic organs. Weakening from childbirth, age, or chronic straining allows the bladder, rectum, uterus, or vaginal apex to descend. The dropping of an organ is graded I-IV based on descent relative to the hymen. Symptoms result from the bulging mass and associated organ dysfunction.

\section*{Causes}
\begin{itemize}
\item Vaginal childbirth (most common risk factor)
\item Chronic constipation and straining
\item Chronic cough (COPD, smoking)
\item Obesity
\item Heavy lifting
\item Aging and menopause (estrogen loss weakens tissues)
\item Connective tissue disorders
\item Prior pelvic surgery (hysterectomy)
\end{itemize}

\section*{When to See a Doctor}
\begin{itemize}
\item Sensation of a bulge or pressure in the vagina
\item Feeling of something falling out of the vagina
\item Lower back pain or pelvic pressure
\item Urinary symptoms: incontinence, frequency, urgency, incomplete emptying
\item Difficulty with bowel movements (need to splint the vagina)
\item Sexual discomfort
\item Vaginal discharge or bleeding
\item Symptoms worsen with standing, lifting, or prolonged activity and improve when lying down
\end{itemize}

\section*{Related Symptoms}
\begin{itemize}
\item Urinary incontinence (stress or urge)
\item Overactive bladder
\item Frequent urination
\item Constipation
\item Rectocele (difficulty with bowel movements)
\item Cystocele (bladder dropping of an organ)
\end{itemize}

\section*{Self-Care / Home Management}
\begin{itemize}
\item Pelvic floor exercises (Kegel exercises) daily
\item Maintain a healthy weight
\item Avoid heavy lifting and straining
\item Treat constipation with high-fiber diet and hydration
\item Use vaginal estrogen cream if postmenopausal
\item Consider a vaginal pessary for support
\item Quit smoking to reduce chronic cough
\end{itemize}

\section*{How Your Doctor Will Evaluate You}
\begin{itemize}
\item Pelvic examination with Valsalva to assess descent
\item Pelvic organ dropping of an organ quantification (POP-Q) staging
\item Urodynamic testing if urinary symptoms are present
\item Ultrasound or MRI for complex cases
\end{itemize}

\section*{Treatment Options}
\begin{itemize}
\item Pelvic floor physical therapy
\item Pessary fitting (vaginal support device)
\item Vaginal estrogen therapy
\item Surgery: sacrocolpopexy, uterosacral ligament suspension, anterior/posterior repair, colpocleisis for severe cases
\end{itemize}
